\begin{solution}{101}
    \begin{subsolution}
        当 $n_{1}>0$ ， $n_{2}>0$ 时，如图 1a 所示，AB 为入射光的等光程（等相位）面，当 B 点到达 $B'$ 时，A 点发出的子波的波前在介质 2 中位于一个以 A 点为球心的半球面上，此球面在光线的入射面上是半径为 $r$ 的半圆，半径 $r$ 由下式确定
        \begin{equation}
            n _ {2} r = n _ {1} \left| B B ^ {\prime} \right|
            \label{eq:1.p101}
        \end{equation}
        做直线 $B^{\prime}A^{\prime}$ 与子波相位面的圆相切于 $A^{\prime}$ 点，则 $B^{\prime}A^{\prime}$ 是等光程（等相位）面。由几何关系
        \begin{align*}
            \left| B B ^ {\prime} \right| = \left| A B ^ {\prime} \right| \sin \theta_ {1} \\
            r = \left| A B ^ {\prime} \right| \sin \theta_ {2}
        \end{align*}
        由以上方程得到
        \begin{equation}
            n _ {1} \sin \theta_ {1} = n _ {2} \sin \theta_ {2}
            \label{eq:2.p101}
        \end{equation}
        当 $n_{1}>0$ ， $n_{2}<0$ 时，在介质 2 中传播的光的光程为负。

        (解法一)

        如图 1b 所示，在 $B'$ 发出的子波的波前是半径为 $r$ 的半圆，对应于光程 $n_{2}r < 0$ 。因此，在介质 2 中与 A 点等相位的面由通过 A 点与圆相切的线确定，切点为 $A'$ 。圆的半径由下式确定
        \begin{equation}
            n _ {1} \left| B B ^ {\prime} \right| + n _ {2} \left| B ^ {\prime} A ^ {\prime} \right| = 0
            \label{eq:3.p101}
        \end{equation}
        由几何关系
        \begin{equation*}
            \left| B B ^ {\prime} \right| = \left| A B ^ {\prime} \right| \sin \theta_ {1}, \quad \left| B ^ {\prime} A ^ {\prime} \right| = \left| A B ^ {\prime} \right| \sin \left| \theta_ {2} \right|
        \end{equation*}
        由以上方程得到
        \begin{equation}
            n _ {1} \left| A B ^ {\prime} \right| \sin \theta_ {1} = - n _ {2} \left| A B ^ {\prime} \right| \sin \left| \theta_ {2} \right| = n _ {2} \left| A B ^ {\prime} \right| \sin \theta_ {2}
            \label{eq:4.p101}
        \end{equation}
        即
        \begin{equation}
            n _ {1} \sin \theta_ {1} = n _ {2} \sin \theta_ {2}
            \label{eq:5.p101}
        \end{equation}
        [（解法二）

        如图 1b' 所示，AB 为入射光的等光程（等相位）面，当 B 点到达 $B'$ 时，A 点发出的子波的波前在介质 2 中的入射面上是以 A 点为圆心，半径为 $r$ 的圆，对应于光程 $n_{2}r < 0$ 。半径 $r$ 由下式确定
        \begin{equation*}
            \left| n _ {2} \right| r = n _ {1} \left| B B ^ {\prime} \right|
        \end{equation*}
        因光程为负，等效于等光程面反向移动，成为在介质 1 中以虚线画出的半圆。做直线 $B'A'$ 与虚线的子波相位面的圆相切于 $A'$ 点，则 $B'A'$ 是等光程（等相位）面。由几何关系
        \begin{equation*}
            \left| B B ^ {\prime} \right| = \left| A B ^ {\prime} \right| \sin \theta_ {1}, r = \left| A B ^ {\prime} \right| \sin \left| \theta_ {2} \right|
        \end{equation*}
        由以上方程得到
        \begin{equation*}
            n _ {1} \sin \theta_ {1} = n _ {2} \sin \theta_ {2}
        \end{equation*}
        ]
    \end{subsolution}
    \begin{subsolution}
        约定物侧的物距 $s_{1}$ 和像侧的像距 $s_{2}$ 取为正，物侧的像距和像侧的物距取为负。入射角 $\theta_{1}$ 为正，折射角 $\theta_{2}$ 为负。其余图中标出的角度均取为正。由图 c 可见
        \begin{equation}
            \theta_ {1} = \alpha_ {1} + \beta , - \theta_ {2} = \alpha_ {2} - \beta
            \label{eq:6.p101}
        \end{equation}
        在傍轴条件下，这些角度均为小角度，有
        \begin{equation*}
            \sin \theta_ {1} = \theta_ {1}, \quad \sin \theta_ {2} = \theta_ {2}
        \end{equation*}
        以及
        \begin{align*}
            \tan \alpha_ {1} = \sin \alpha_ {1} = \alpha_ {1} = \frac {h}{s _ {1}} \\
            \tan \alpha_ {2} = \sin \alpha_ {2} = \alpha_ {2} = \frac {h}{s _ {2}}
        \end{align*}
        \begin{equation}
            \tan \beta = \sin \beta = \beta = \frac {h}{R}
            \label{eq:7.p101}
        \end{equation}
        这里，$h$ 是折射点 M 与光轴的距离。由折射定律
        \begin{equation*}
            n _ {1} \sin \theta_ {1} = n _ {2} \sin \theta_ {2}
        \end{equation*}
        即
        \begin{equation}
            n _ {1} \left(\frac {h}{s _ {1}} + \frac {h}{R}\right) = - n _ {2} \left(\frac {h}{s _ {2}} - \frac {h}{R}\right)
            \label{eq:8.p101}
        \end{equation}
        整理得到
        \begin{equation}
            \frac {n _ {1}}{s _ {1}} + \frac {n _ {2}}{s _ {2}} = \frac {n _ {2} - n _ {1}}{R}
            \label{eq:9.p101}
        \end{equation}
        注意到 $n_{2}<0,\quad n_{1}>0$ ，上式右方为负。

        为计算横向放大率，如图 2 所示，考虑物 $Q_{1}P_{1}$ ，设像为 $Q_{2}P_{2}$ 。设想光轴绕 C 点转过一个小角，使得 $P_{1}$ 点位于光轴上，原 $Q_{1}$ 转到 $Q_{1}^{\prime}$ ， $Q_{2}$ 转到 $Q_{2}^{\prime}$ ， $P_{1}$ 的像为 $P_{2}$ ，在傍轴近似下，因转角很小
        \begin{align}
            Q _ {1} P _ {1} = Q _ {1} Q _ {1} ^ {\prime} \\
            Q _ {2} P _ {2} = Q _ {2} Q _ {2} ^ {\prime}
            \label{eq:10.p101}
        \end{align}
        由图 2 的几何关系有
        \begin{align}
            \left| Q _ {1} P _ {1} \right| = s _ {1} \theta_ {1} \\
            \left| Q _ {2} P _ {2} \right| = - s _ {2} \theta_ {2}
            \label{eq:12.p101}
        \end{align}
        由此得到横向放大率公式
        \begin{equation}
            \frac {\left| Q _ {2} P _ {2} \right|}{\left| Q _ {1} P _ {1} \right|} = \frac {- s _ {2} \theta_ {2}}{s _ {1} \theta_ {1}} = - \frac {s _ {2} n _ {1}}{s _ {1} n _ {2}}
            \label{eq:13.p101}
        \end{equation}
    \end{subsolution}
    \begin{subsolution}
        如图 3 所示, 如果没有球面隔开的介质 2, 平行光经过薄凸透镜后将汇聚于焦点，如图中虚线所示。在引入介质 2 后，虚线所示的汇聚点是球面的虚物，对应的物距是
        \begin{equation}
            s _ {1} = - (f - d)
            \label{eq:14.p101}
        \end{equation}
        如果 $n_2 > 0$ ，已知球面的成像公式为
        \begin{equation}
            \frac {n _ {1}}{s _ {1}} + \frac {n _ {2}}{s _ {2}} = \frac {n _ {2} - n _ {1}}{R}
            \label{eq:15.p101}
        \end{equation}
        形式上与（2）中求出的成像公式完全相同，在 $n_2 > 0$ 和 $n_2 < 0$ 的情况下像距 $s_2$ 都可以解出为
        \begin{equation}
            s _ {2} = \frac {n _ {2} s _ {1} R}{n _ {2} s _ {1} - n _ {1} (s _ {1} + R)} = \frac {n _ {2} R (f - d)}{(n _ {2} - n _ {1}) (f - d) + n _ {1} R}
            \label{eq:16.p101}
        \end{equation}
        对于题中给出的数据，求得各自的 $s_2$ ，如下表所示：

        \begin{tabular}{|c|c|c|c|}
            \hline
            序号 & $n_2$ & $d$ & $s_2$ \\
            \hline
            1 & 1.5 & 0.35R & 1.10R \\
            \hline
            2 & 1.5 & 0.85R & 0.74R \\
            \hline
            3 & -1.5 & 0.35R & 0.92R \\
            \hline
            4 & -1.5 & 0.85R & 1.56R \\
            \hline
        \end{tabular}

        按照表中结果，画出的示意图由图 3 给出。
    \end{subsolution}
\end{solution}